\documentclass[10pt,twocolumn,a4paper]{article}
\usepackage[margin=1.8cm]{geometry}
\usepackage{setspace}
\onehalfspacing
\usepackage{times}
\usepackage{amsmath,amssymb,amsfonts}
\usepackage{graphicx}
\usepackage{booktabs}
\usepackage{multirow}
\usepackage{url}
\usepackage{float}
\usepackage{cite}
\usepackage{xcolor}
\usepackage{caption}
\usepackage{titlesec}
\usepackage{fancyhdr}
\usepackage[utf8]{inputenc}
\usepackage[T1]{fontenc}
% babel turkish disabled for Overleaf compatibility

\captionsetup{font=small,labelfont=bf}
\titleformat{\section}{\normalfont\large\bfseries\centering}{\thesection.}{0.5em}{}
\titleformat{\subsection}{\normalfont\normalsize\bfseries}{\thesubsection}{0.5em}{}
\pagestyle{fancy}
\fancyhf{}
\fancyhead[L]{\small Tar{\i}mda Ak{\i}ll{\i} G\"ubreleme ve Sulama i\c{c}in DQN}
\fancyhead[R]{\small IEEE Format}
\fancyfoot[C]{\thepage}
\renewcommand{\headrulewidth}{0.4pt}

\begin{document}

\title{\vspace{-0.6em}\textbf{Tar{\i}mda Ak{\i}ll{\i} G\"ubreleme ve Sulama i\c{c}in Derin Q-A\u{g}{\i} Tabanl{\i} Yakla\c{s}{\i}m}}

\author{
Metehan S\"ozenli\\
\textit{Bilgisayar M\"uhendisli\u{g}i B\"ol\"um\"u, Y\"uksek Lisans}\\
Bursa Teknik \"Universitesi, Bursa, T\"urkiye\\
\texttt{25435004037@ogr.btu.edu.tr}\\[0.35em]
Do\c{c}.~Dr.~Kubilay Demir\\
\textit{Dan{\i}\c{s}man --- Derin Peki\c{s}tirmeli \"O\u{g}renme Dersi}
}

\date{}
\maketitle
\thispagestyle{fancy}

\begin{abstract}
\noindent
Sulama ve g\"ubreleme kararlar{\i} s{\i}ral{\i} bir karar verme problemi olarak modellenebilir. Bu \c{c}al{\i}\c{s}mada, bitki verimini art{\i}rmay{\i} ve a\c{s}{\i}r{\i} g\"ubre kaynakl{\i} toksisiteyi s{\i}n{\i}rlamay{\i} hedefleyen bir Derin Q-A\u{g}{\i} (DQN) uygulamas{\i} sunulmaktad{\i}r. Ortam, ger\c{c}ek bir \"ur\"un verim veri setinden al{\i}nan istatistikler ile \c{c}arpanl{\i} stres tabanl{\i} basit bir b\"uy\"ume modelinin birle\c{s}tirildi\u{g}i hibrit bir Markov Karar S\"ureci olarak kurulmu\c{s}tur. Durum; nem, besin, sa\u{g}l{\i}k ve toksisite gibi s\"urekli de\u{g}i\c{s}kenleri i\c{c}erir. Ajan sekiz ayr{\i}k aksiyon aras{\i}ndan se\c{c}im yapar; g\"ubre dozlar{\i} veri da\u{g}{\i}l{\i}m{\i}na yakla\c{s}{\i}k hizalanm{\i}\c{s}t{\i}r. Deneylerde e\u{g}itilmi\c{s} politika, rastgele politikaya g\"ore daha y\"uksek verim ve daha d\"u\c{s}\"uk g\"ubre kullan{\i}m{\i} \"uretmi\c{s}tir. Sonu\c{c}lar, ideal g\"unl\"uk tarla zaman serisinin yoklu\u{g}u ve problemin g\"oreli basitli\u{g}i g\"oz \"on\"unde bulundurularak yorumlanmal{\i}d{\i}r.
\end{abstract}

\begin{center}
\small\textbf{Anahtar Kelimeler---}Derin Peki\c{s}tirmeli \"O\u{g}renme, DQN, Hassas Tar{\i}m, G\"ubreleme Optimizasyonu, Hibrit Sim\"ulasyon, Verim Y\"onetimi
\end{center}

\section{Giri\c{s}}
Sulama ve kimyasal g\"ubre y\"onetimi modern tar{\i}m{\i}n en kritik operasyonel kararlar{\i}ndan birini olu\c{s}turmaktad{\i}r. A\c{s}{\i}r{\i} g\"ubreleme \"uretim maliyetini y\"ukseltmekte, toprak yap{\i}s{\i}n{\i} bozmakta, y\"uzey ve yeralt{\i} sular{\i}n{\i} kirletmektedir. Yetersiz uygulama ise verim kayb{\i}na yol a\c{c}maktad{\i}r~\cite{tilman2002agricultural}. Bu kararlar do\u{g}as{\i} gere\u{g}i s{\i}ral{\i}d{\i}r: her g\"unl\"uk m\"udahale stokastik hava ko\c{s}ullar{\i} alt{\i}nda toprak nemini, besin durumunu, bitki sa\u{g}l{\i}\u{g}{\i}n{\i} ve sonraki g\"unlerin b\"uy\"ume potansiyelini etkilemektedir.

Peki\c{s}tirmeli \"o\u{g}renme (RL) s{\i}ral{\i} karar verme i\c{c}in bi\c{c}imsel bir \c{c}er\c{c}eve sunmaktad{\i}r. Tabular Q-\"O\u{g}renme y\"ontemleri d\"u\c{s}\"uk boyutlu ayr{\i}k ortamlarda uygulanabilmi\c{s} olsa da s\"urekli sens\"or okumalar{\i} (nem, besin konsantrasyonu, s{\i}cakl{\i}k, pH vb.) devreye girdi\u{g}inde boyutlanabilirlik sorunu ortaya \c{c}{\i}kmaktad{\i}r. Derin Q-A\u{g}lar{\i} (DQN) eylem-de\u{g}er fonksiyonunu sinir a\u{g}{\i} ile yakla\c{s}{\i}kla\c{s}t{\i}rarak, deneyim tekrar{\i} ve hedef a\u{g} ile e\u{g}itimi stabilize ederek bu s{\i}n{\i}rlamay{\i} a\c{s}maya y\"oneliktir~\cite{mnih2015human}.

Bu \c{c}al{\i}\c{s}mada tabular form\"ulasyon yerine tam bir DQN ajan{\i} kullan{\i}lmakta ve ortam ger\c{c}ek bir \"ur\"un verim veri seti \"uzerinde kalibre edilmektedir. Taban verim istatistikleri ve g\"ubre miktar da\u{g}{\i}l{\i}mlar{\i} ger\c{c}ek veriden al{\i}nm{\i}\c{s}t{\i}r. G\"unl\"uk ge\c{c}i\c{s}ler ise FAO AquaCrop~\cite{raes2009aquacrop} ve Doorenbos ile Kassam'{\i}n klasik verim-tepki fakt\"or\"unden~\cite{doorenbos1979} esinlenen hafif fizik tabanl{\i} bir b\"uy\"ume modeli ile y\"onetilmektedir. Sistem toksisiteyi ve kaynak israf{\i}n{\i} cezaland{\i}r{\i}rken verimi art{\i}ran politikalar \"o\u{g}renmektedir.

\c{C}al{\i}\c{s}mada \c{s}u ad{\i}mlar at{\i}lm{\i}\c{s}t{\i}r: (i) ger\c{c}ek verim ve g\"ubre istatistikleri ile kalibre edilmi\c{s} Gymnasium uyumlu hibrit bir tar{\i}m ortam{\i}, (ii) deneyim tekrar{\i} ve hedef a\u{g} i\c{c}eren tam bir PyTorch DQN uygulamas{\i}, (iii) \"od\"ul, verim, zamansal fark kayb{\i}, ke\c{s}if oran{\i} ve aksiyon da\u{g}{\i}l{\i}m{\i}n{\i} kapsayan kapsaml{\i} deneysel de\u{g}erlendirme, (iv) \"o\u{g}renilen politikan{\i}n g\"orsel sim\"ulasyonu.

\section{Literat\"ur Taramas{\i}}
\subsection{Tar{\i}mda Peki\c{s}tirmeli \"O\u{g}renme}
Tar{\i}mda peki\c{s}tirmeli \"o\u{g}renme uygulamalar{\i} son on y{\i}lda belirgin bir art{\i}\c{s} g\"ostermi\c{s}tir. Erken \c{c}al{\i}\c{s}malar tabular Q-\"O\u{g}renme ve SARSA ile sulama zamanlamas{\i} ve g\"ubre \"onerisi problemlerini ele alm{\i}\c{s}t{\i}r. Bu y\"ontemler d\"u\c{s}\"uk boyutlu ayr{\i}k ortamlarda i\c{s}e yarasa da s\"urekli sens\"or de\u{g}i\c{s}kenleri eklendi\u{g}inde durum uzay{\i} h{\i}zla b\"uy\"umekte ve tabular yakla\c{s}{\i}mlar pratik olmaktan \c{c}{\i}kmaktad{\i}r. Derin RL y\"ontemleri bu sorunu a\c{s}mak i\c{c}in devreye girmi\c{s}tir.

Mnih ve arkada\c{s}lar{\i}n{\i}n DQN algoritmas{\i}~\cite{mnih2015human} deneyim tekrar{\i} ve hedef a\u{g} ile derin sinir a\u{g}lar{\i}n{\i} Q-\"o\u{g}renmeye entegre ederek Atari oyunlar{\i}nda insan d\"uzeyinde performans g\"ostermi\c{s}tir. Bu mimari s\"urekli veya y\"uksek boyutlu durum uzaylar{\i}na sahip tar{\i}msal kontrol problemlerine do\u{g}rudan uyarlanabilmektedir. Sonraki iyile\c{s}tirmeler aras{\i}nda Double DQN, Dueling DQN ve Prioritized Experience Replay yer almaktad{\i}r. Akt\"or-ele\c{s}tirmen ailesinden A2C, PPO ve SAC algoritmalar{\i} ise s\"urekli aksiyon uzaylar{\i} i\c{c}in tercih edilmektedir~\cite{wang2025fertilization}.

CropGym~\cite{overweg2021cropgym} PCSE s\"ure\c{c} tabanl{\i} bitki sim\"ulat\"or\"u ile RL ajan{\i}n{\i} birle\c{s}tiren a\c{c}{\i}k kaynak bir Gymnasium ortam{\i}d{\i}r. Ajan g\"unl\"uk azot uygulamas{\i} kararlar{\i} vermekte; \"od\"ul verim ile g\"ubre maliyeti aras{\i}ndaki dengeyi yans{\i}tmaktad{\i}r. SWAT tabanl{\i} RL \c{c}er\c{c}eveleri~\cite{madondo2023swat} ve DSSAT entegrasyonlar{\i} da raporlanm{\i}\c{s}t{\i}r. Bu \c{c}al{\i}\c{s}malar y\"uksek do\u{g}ruluklu fiziksel sim\"ulat\"orlerin RL i\c{c}in uygun ancak hesaplama a\c{c}{\i}s{\i}ndan pahal{\i} oldu\u{g}unu g\"ostermektedir. Buna kar\c{s}{\i}l{\i}k bu \c{c}al{\i}\c{s}mada binlerce b\"ol\"um\"un h{\i}zl{\i} \c{c}al{\i}\c{s}t{\i}r{\i}labilmesi i\c{c}in basitle\c{s}tirilmi\c{s} bir ge\c{c}i\c{s} modeli kullan{\i}l{\i}rken temel parametreler ger\c{c}ek veri da\u{g}{\i}l{\i}mlar{\i}na ba\u{g}lanmaktad{\i}r.

B\"u et ve arkada\c{s}lar{\i}~\cite{buet2022deep} derin RL'nin sulama planlamas{\i}nda klasik optimizasyon y\"ontemlerine g\"ore daha esnek oldu\u{g}unu g\"ostermi\c{s}tir. Gautron ve arkada\c{s}lar{\i}~\cite{gautron2022gym} tar{\i}m i\c{c}in standart RL ortamlar{\i}n{\i}n eksikli\u{g}ine dikkat \c{c}ekerek Gym uyumlu aray\"uzlerin \"onemini vurgulam{\i}\c{s}t{\i}r. Bu \c{c}al{\i}\c{s}ma da Gymnasium uyumlu bir ortam sunarak tekrarlanabilirli\u{g}i desteklemektedir.

\subsection{Verim-Stres \.{I}li\c{s}kileri ve Fiziksel Modeller}
Su stresi ile g\"oreli verim aras{\i}ndaki ili\c{s}ki FAO verim-tepki fakt\"or\"u $K_y$ ile tan{\i}mlanmaktad{\i}r~\cite{doorenbos1979}:
\begin{equation}
\left(1-\frac{Y_a}{Y_m}\right)=K_y\left(1-\frac{ET_a}{ET_m}\right).
\end{equation}
Burada $Y_a$ ve $Y_m$ s{\i}ras{\i}yla ger\c{c}ek ve maksimum verim, $ET_a$ ve $ET_m$ ger\c{c}ek ve maksimum evapotranspirasyonu ifade etmektedir. $K_y>1$ bitkinin su a\c{c}{\i}\u{g}{\i}na duyarl{\i} oldu\u{g}unu, $K_y<1$ ise daha toleransl{\i} oldu\u{g}unu g\"ostermektedir. Bu form\"ulasyon bir\c{c}ok tah{\i}l ve sebze t\"ur\"u i\c{c}in deneysel olarak do\u{g}rulanm{\i}\c{s}t{\i}r.

AquaCrop bu fikri genelle\c{s}tirerek stres katsay{\i}lar{\i}n{\i} kanopi geni\c{s}lemesi, stoma iletkenli\u{g}i, ya\c{s}lanma ve hasat indeksine ayr{\i} ayr{\i} uygulamaktad{\i}r~\cite{raes2009aquacrop,hsiao2009aquacrop,steduto2009aquacrop}. Model su verimlili\u{g}i tabanl{\i} bir biyok\"utle \"uretimi form\"ulasyonu kullanmakta ve bir\c{c}ok \"ur\"un i\c{c}in parametrelenmi\c{s}tir. DSSAT ve APSIM gibi di\u{g}er s\"ure\c{c} tabanl{\i} modeller de benzer stres mekanizmalar{\i}n{\i} i\c{c}ermektedir~\cite{jones2003dssat}. Bu \c{c}al{\i}\c{s}mada kullan{\i}lan b\"uy\"ume modeli AquaCrop'un \c{c}arpanl{\i} stres yap{\i}s{\i}n{\i} sadele\c{s}tirilmi\c{s} haliyle benimsemektedir; b\"oylece tar{\i}msal yorumlanabilirlik korunurken RL i\c{c}in hesaplama maliyeti d\"u\c{s}\"uk tutulmaktad{\i}r.

Azotlu g\"ubrelerin a\c{s}{\i}r{\i} kullan{\i}m{\i} toksisite, tuzluluk art{\i}\c{s}{\i} ve verim kayb{\i}na yol a\c{c}t{\i}\u{g}{\i} literat\"urde belgelenmi\c{s}tir~\cite{salvagiotti2009nitrogen,cameron2013nitrogen}. Azot d\"ong\"us\"undeki kay{\i}plar hem \c{c}evresel hem ekonomik maliyet yaratmaktad{\i}r. Bu nedenle \"od\"ul fonksiyonunda toksisite cezas{\i} a\c{c}{\i}k\c{c}a yer almaktad{\i}r. Besin eksikli\u{g}i ve fazlal{\i}\u{g}{\i}n{\i}n her ikisinin de stresi art{\i}rd{\i}\u{g}{\i} bilinmektedir; modelde besin stresi hem d\"u\c{s}\"uk hem y\"uksek seviyelerde cezaland{\i}r{\i}lmaktad{\i}r.

\subsection{Hassas Tar{\i}m ve Veri Odakl{\i} Yakla\c{s}{\i}mlar}
Hassas tar{\i}m mek\^ansal ve zamansal de\u{g}i\c{s}kenli\u{g}i dikkate alarak girdi kullan{\i}m{\i}n{\i} optimize etmeyi ama\c{c}lamaktad{\i}r~\cite{zhang2002precision}. IoT toprak sens\"orleri, hava istasyonlar{\i}, uydu ve drone g\"or\"unt\"uleri ile makine \"o\u{g}renmesi teknikleri bu alanda yayg{\i}nla\c{s}m{\i}\c{s}t{\i}r. Denetimli \"o\u{g}renme y\"ontemleri verim tahmini i\c{c}in s{\i}k\c{c}a kullan{\i}lsa da s{\i}ral{\i} karar verme problemi i\c{c}in RL daha do\u{g}al bir \c{c}er\c{c}eve sunmaktad{\i}r \c{c}\"unk\"u gelecek \"od\"ullerin iskonto edilmi\c{s} toplam{\i}n{\i} maksimize etmektedir~\cite{kamilaris2018deep}.

Ger\c{c}ek tar{\i}msal veri setlerinin g\"unl\"uk RL ortamlar{\i}na aktar{\i}lmas{\i} veri k{\i}tl{\i}\u{g}{\i} nedeniyle genellikle kalibrasyon yoluyla yap{\i}lmaktad{\i}r. Bu \c{c}al{\i}\c{s}mada da ayn{\i} strateji izlenmektedir: taban verim ve g\"ubre istatistikleri ger\c{c}ek veriden al{\i}nmakta, g\"unl\"uk dinamikler ise literat\"ure dayal{\i} stres modeli ile \"uretilmektedir. Bu yakla\c{s}{\i}m tamamen sentetik ortamlara k{\i}yasla daha y\"uksek d{\i}\c{s} ge\c{c}erlilik sa\u{g}lamaktad{\i}r~\cite{overweg2021cropgym}.

\section{Problem Form\"ulasyonu}
Otuz g\"unl\"uk \"ur\"un d\"ong\"us\"u sonlu ufuklu bir Markov Karar S\"ureci (MDP) olarak modellenmi\c{s}tir. MDP $(S,A,P,R,\gamma)$ be\c{s}lisi ile tan{\i}mlanmaktad{\i}r; burada $S$ durum uzay{\i}n{\i}, $A$ aksiyon uzay{\i}n{\i}, $P$ ge\c{c}i\c{s} olas{\i}l{\i}klar{\i}n{\i}, $R$ \"od\"ul fonksiyonunu ve $\gamma$ indirim fakt\"or\"un\"u ifade etmektedir.

\subsection{Durum Uzay{\i}}
Her g\"un $t$ ajan \c{s}u s\"urekli vekt\"or\"u g\"ozlemlemektedir:
\begin{equation}
s_t = [\tilde{t},\, m_t,\, n_t,\, pH_t,\, T_t,\, H_t,\, L_t,\, h_t,\, \tau_t] \in \mathbb{R}^9.
\end{equation}
Burada $\tilde{t}=t/30$ normalize g\"un indeksini, $m_t$ toprak nemini (\%), $n_t$ g\"oreli besin seviyesini, $pH_t$ asitli\u{g}i, $T_t$ s{\i}cakl{\i}\u{g}{\i}, $H_t$ ba\u{g}{\i}l nemi, $L_t$ {\i}\c{s}{\i}k \c{s}iddetini, $h_t$ bitki sa\u{g}l{\i}k skorunu (0--100) ve $\tau_t$ toksisite indeksini g\"ostermektedir. Bu de\u{g}i\c{s}kenler tipik IoT toprak ve hava sens\"orlerinden elde edilebilecek \"ol\c{c}\"umlerle \"ort\"u\c{s}mektedir.

\subsection{Aksiyon Uzay{\i}}
Sekiz ayr{\i}k aksiyon ger\c{c}ek\c{c}i y\"onetim kararlar{\i}n{\i} kodlamaktad{\i}r (Tablo~\ref{tab:actions}). G\"ubre dozlar{\i} ger\c{c}ek veri setinde g\"ozlenen ampirik da\u{g}{\i}l{\i}m{\i}n i\c{c}inde yer almaktad{\i}r (medyan yakla\c{s}{\i}k 137\,kg/ha). Bu sayede ajan{\i}n se\c{c}ebilece\u{g}i g\"ubre miktarlar{\i} ger\c{c}ek tar{\i}msal uygulamalarla uyumludur. Tek uygulama dozlar{\i} 40, 70 ve 140\,kg/ha olarak belirlenmi\c{s}; b\"oylece iki ile d\"ort uygulama ile sezonluk toplam{\i}n ger\c{c}ek\c{c}i aral{\i}\u{g}a ula\c{s}mas{\i} m\"umk\"un k{\i}l{\i}nm{\i}\c{s}t{\i}r.

\begin{table}[H]
\centering
\caption{Ayr{\i}k Aksiyon K\"umesi (ger\c{c}ek veri y\"uzdeliklerine hizal{\i})}
\label{tab:actions}
\small
\begin{tabular}{@{}clcc@{}}
\toprule
\.{I}ndeks & Ad & Su (mm) & G\"ubre (kg/ha) \\
\midrule
0 & Bekle & 0 & 0 \\
1 & Hafif Sulama & 14 & 0 \\
2 & Yo\u{g}un Sulama & 28 & 0 \\
3 & G\"ubre 40\,kg & 0 & 40 \\
4 & G\"ubre 70\,kg & 0 & 70 \\
5 & G\"ubre 140\,kg & 0 & 140 \\
6 & Sulama + G\"ubre 50\,kg & 16 & 50 \\
7 & Sulama + G\"ubre 90\,kg & 18 & 90 \\
\bottomrule
\end{tabular}
\end{table}

\subsection{Ge\c{c}i\c{s} Dinamikleri ve B\"uy\"ume Modeli}
Aksiyon uyguland{\i}ktan sonra toprak de\u{g}i\c{s}kenleri (nem, besin, toksisite) g\"uncellenmekte; ard{\i}ndan su, besin ve toksisite stresleri hesaplanmaktad{\i}r. G\"unl\"uk g\"oreli b\"uy\"ume literat\"urdeki \c{c}arpanl{\i} stres yakla\c{s}{\i}m{\i}na uygun olarak \c{s}u \c{s}ekilde ifade edilmektedir~\cite{raes2009aquacrop,doorenbos1979}:
\begin{equation}
g_t = g_0 \cdot (1-k_w\sigma_w)\cdot(1-k_n\sigma_n)\cdot(1-k_\tau\sigma_\tau).
\end{equation}
Burada $g_0$ potansiyel g\"unl\"uk b\"uy\"umeyi, $\sigma_w$, $\sigma_n$ ve $\sigma_\tau$ s{\i}ras{\i}yla su, besin ve toksisite streslerini, $k_w$, $k_n$ ve $k_\tau$ ise stres katsay{\i}lar{\i}n{\i} g\"ostermektedir. K\"um\"ulatif b\"uy\"ume $G$ ve toksisite cezas{\i} nihai verimi belirlemektedir:
\begin{equation}
Y = Y_{\mathrm{base}} \cdot G \cdot \bigl(1-\min(\tau/25,\,0.55)\bigr).
\end{equation}
Taban verim $Y_{\mathrm{base}}$ ger\c{c}ek veri setinin ampirik da\u{g}{\i}l{\i}m{\i}ndan \"orneklenmektedir; b\"oylece model \c{c}{\i}kt{\i}lar{\i} ger\c{c}ek verim aral{\i}\u{g}{\i} ile tutarl{\i} kalmaktad{\i}r.

\subsection{\"Od\"ul Fonksiyonu}
\c{S}ekillendirilmi\c{s} \"od\"ul sa\u{g}l{\i}kl{\i} b\"uy\"umeyi te\c{s}vik etmekte, toksisiteyi ve kaynak harcamas{\i}n{\i} cezaland{\i}rmakta, d\"onem sonunda verim bonusu vermektedir. Ayr{\i}ca sezonluk toplam g\"ubre kullan{\i}m{\i}n{\i}n ger\c{c}ek veri medyan{\i}na yak{\i}n olmas{\i} \"od\"ullendirilmektedir. \.{I}ndirim fakt\"or\"u $\gamma=0.97$ olarak se\c{c}ilmi\c{s}tir. \"Od\"ul tasar{\i}m{\i} hem k{\i}sa vadeli sa\u{g}l{\i}k iyile\c{s}tirmelerini hem de uzun vadeli hasat hedefini dengeleyecek \c{s}ekilde ayarlanm{\i}\c{s}t{\i}r.

\section{DQN Mimarisi ve E\u{g}itim}
\subsection{A\u{g} Mimarisi}
Eylem-de\u{g}er fonksiyonu $Q(s,a;\theta)$ \c{c}ok katmanl{\i} alg{\i}lay{\i}c{\i} (MLP) ile yakla\c{s}{\i}kla\c{s}t{\i}r{\i}lmaktad{\i}r. Giri\c{s} katman{\i} 9 boyutlu durumu almakta; 256, 128 ve 64 boyutlu \"u\c{c} gizli katman ReLU aktivasyonu ve Katman Normalizasyonu kullanmakta; \c{c}{\i}k{\i}\c{s} katman{\i} her ayr{\i}k aksiyon i\c{c}in bir Q-de\u{g}eri \"uretmektedir. Katman Normalizasyonu e\u{g}itim kararl{\i}l{\i}\u{g}{\i}n{\i} art{\i}rmak i\c{c}in eklenmi\c{s}tir~\cite{ba2016layer}.

\subsection{E\u{g}itim Algoritmas{\i}}
E\u{g}itim klasik DQN algoritmas{\i}n{\i}~\cite{mnih2015human} takip etmektedir:
\begin{itemize}
    \item Deneyim tekrar{\i} tamponu kapasitesi 60\,000 ge\c{c}i\c{s}tir.
    \item Mini-y{\i}\u{g}{\i}n boyutu 128'dir.
    \item Hedef a\u{g} her 300 gradyan ad{\i}m{\i}nda sert g\"uncellenmektedir.
    \item $\varepsilon$-a\c{c}g\"ozl\"u ke\c{s}if 1.0'dan 0.05'e \"ustel azalmaktad{\i}r.
    \item Optimizasyon i\c{c}in Adam algoritmas{\i} kullan{\i}lmakta, \"o\u{g}renme oran{\i} $1\times10^{-4}$ olarak ayarlanmaktad{\i}r.
    \item Kay{\i}p fonksiyonu Smooth L1 (Huber) olarak se\c{c}ilmi\c{s}tir.
    \item Gradyan normu 5.0 ile k{\i}rp{\i}lmaktad{\i}r.
\end{itemize}
Bir b\"ol\"um 30 g\"unl\"uk ad{\i}mdan olu\c{s}maktad{\i}r. Toplam 650 b\"ol\"um (yakla\c{s}{\i}k 19\,500 ortam etkile\c{s}imi) e\u{g}itilmi\c{s}tir. Tekrarlanabilirlik i\c{c}in rastgele tohumlar sabitlenmi\c{s}tir.

\subsection{Ger\c{c}ek Veri Entegrasyonu}
Bu \c{c}al{\i}\c{s}man{\i}n temel \"ozelliklerinden biri ortam{\i}n ger\c{c}ek tar{\i}msal veri ile kalibre edilmesidir. Her b\"ol\"um\"un ba\c{s}{\i}nda Hindistan \"ur\"un verim veri setinden bir \"ornek \c{c}ekilmektedir. Bu \"ornek taban verimi $Y_{\mathrm{base}}$ sa\u{g}lamakta ve g\"ubre yo\u{g}unlu\u{g}unun \"oncel da\u{g}{\i}l{\i}m{\i}n{\i} bilgilendirmektedir. G\"unl\"uk ge\c{c}i\c{s}ler stokastik kalmakta, ancak ortaya \c{c}{\i}kan verim da\u{g}{\i}l{\i}m{\i} ampirik veriyle tutarl{\i} olacak \c{s}ekilde k{\i}s{\i}tlanmaktad{\i}r. B\"oylece ajan yaln{\i}zca sentetik varsay{\i}mlara de\u{g}il, ger\c{c}ek tar{\i}msal verim ve g\"ubre istatistikleri \"uzerinde kalibre edilmi\c{s} bir ortamda \"o\u{g}renmektedir. Bu yakla\c{s}{\i}m sim-to-real bo\c{s}lu\u{g}unu azaltmaya y\"onelik pratik bir ad{\i}m olu\c{s}turmaktad{\i}r.

Veri setinde g\"ubre miktar{\i}n{\i}n medyan{\i} yakla\c{s}{\i}k 137\,kg/ha, birinci \c{c}eyrek yakla\c{s}{\i}k 74\,kg/ha ve \"u\c{c}\"unc\"u \c{c}eyrek yakla\c{s}{\i}k 259\,kg/ha'd{\i}r. Aksiyon setindeki dozlar bu y\"uzdeliklere g\"ore belirlenmi\c{s}; \"od\"ul fonksiyonunda sezonluk toplam{\i}n $[70,260]$ aral{\i}\u{g}{\i}nda olmas{\i} ayr{\i}ca \"od\"ullendirilmi\c{s}tir.

\section{Deneysel Sonu\c{c}lar}
\subsection{\"O\u{g}renme E\u{g}rileri}
\c{S}ekil~\ref{fig:full} ba\c{s}l{\i}ca e\u{g}itim tan{\i}lar{\i}n{\i} bir arada g\"ostermektedir.

\begin{figure*}[t]
\centering
\includegraphics[width=0.93\textwidth]{full_dqn_paper_figures.png}
\caption{650 b\"ol\"umden elde edilen DQN e\u{g}itim tan{\i}lar{\i}. (a) B\"ol\"um k\"um\"ulatif \"od\"ul\"u ve 30 b\"ol\"um hareketli ortalamas{\i}. (b) Nihai verim y\"or\"ungesi; kesikli \c{c}izgi ger\c{c}ek veri medyan{\i}n{\i} g\"ostermektedir. (c) Zamansal fark kayb{\i} (logaritmik \"ol\c{c}ek). (d) Ke\c{s}if oran{\i} $\varepsilon$. (e) E\u{g}itimin son \%30'unda aksiyon frekanslar{\i}. (f) E\u{g}itimin ikinci yar{\i}s{\i}nda ger\c{c}ekle\c{s}en verim histogram{\i}.}
\label{fig:full}
\end{figure*}

\"One \c{c}{\i}kan g\"ozlemler \c{s}unlard{\i}r:
\begin{itemize}
    \item K\"um\"ulatif \"od\"ul rastgele politikaya \"ozg\"u negatif veya d\"u\c{s}\"uk de\u{g}erlerden yakla\c{s}{\i}k 300 b\"ol\"um sonra kararl{\i} pozitif b\"olgeye y\"ukselmektedir.
    \item Ortalama nihai verim yakla\c{s}{\i}k 3500\,kg/ha d\"uzeyine ula\c{s}makta; sezonluk toplam g\"ubre kullan{\i}m{\i} ger\c{c}ek veri medyan{\i}na ($\approx$137\,kg/ha) yak{\i}n seyretmektedir.
    \item Zamansal fark kayb{\i} istikrarl{\i} bi\c{c}imde azalmakta; bu da eylem-de\u{g}er fonksiyonunun ba\c{s}ar{\i}l{\i} yakla\c{s}{\i}kla\c{s}t{\i}r{\i}ld{\i}\u{g}{\i}n{\i} g\"ostermektedir.
    \item Ke\c{s}if oran{\i} sorunsuz bi\c{c}imde $\varepsilon=0.05$ taban{\i}na inmektedir.
    \item Ge\c{c} e\u{g}itim aksiyon istatistikleri hafif sulama ve 40\,kg g\"ubre uygulamalar{\i}n{\i} \"one \c{c}{\i}karmakta; y\"uksek tek seferlik dozlar neredeyse hi\c{c} se\c{c}ilmemektedir.
\end{itemize}

\subsection{Nicel De\u{g}erlendirme}
Tablo~\ref{tab:eval} y\"uz de\u{g}erlendirme b\"ol\"um\"u \"uzerinde a\c{c}g\"ozl\"u politika ($\varepsilon=0$) performans{\i}n{\i} \"ozetlemektedir.

\begin{table}[H]
\centering
\caption{A\c{c}g\"ozl\"u Politika De\u{g}erlendirmesi (100 b\"ol\"um)}
\label{tab:eval}
\small
\begin{tabular}{@{}lccc@{}}
\toprule
Politika & Ort.\ Verim (kg/ha) & Std.\ & Ort.\ \"Od\"ul \\
\midrule
Rastgele & 2700 & 450 & $-1200$ \\
Erken DQN (b\"ol\"um 50) & 3200 & 380 & $+200$ \\
Son DQN & \textbf{3520} & 820 & \textbf{$+660$} \\
\bottomrule
\end{tabular}
\end{table}

E\u{g}itilmi\c{s} ajan rastgele tabana g\"ore rastgele tabana g\"ore daha y\"uksek verim g\"ostermektedir. \"O\u{g}renilen politika sezon boyunca hafif sulama ile nemi korumakta ve birka\c{c} kez 40\,kg g\"ubre uygulayarak toplam g\"ubre miktar{\i}n{\i} ger\c{c}ek veri medyan{\i}na yakla\c{s}t{\i}rmaktad{\i}r. Y\"uksek tek seferlik dozlar toksisite nedeniyle tercih edilmemektedir.

\subsection{Politika Davran{\i}\c{s}{\i}}
\"O\u{g}renilen politika incelendi\u{g}inde ajan{\i}n nemi korumak i\c{c}in a\u{g}{\i}rl{\i}kl{\i} olarak hafif sulama se\c{c}ti\u{g}i, besin d\"u\c{s}t\"u\u{g}\"unde 40\,kg g\"ubre uygulad{\i}\u{g}{\i} ve y\"uksek tek seferlik dozlar{\i} (70--140\,kg) neredeyse hi\c{c} kullanmad{\i}\u{g}{\i} g\"or\"ulmektedir. Sezonluk toplam g\"ubre kullan{\i}m{\i} ortalama $\approx$141\,kg/ha olup ger\c{c}ek veri medyan{\i}na (137\,kg/ha) \c{c}ok yak{\i}nd{\i}r. Bu davran{\i}\c{s} \"od\"uldeki toksisite cezas{\i} ile ger\c{c}ek\c{c}i g\"ubre aral{\i}\u{g}{\i} bonusunun ortak sonucudur ve tar{\i}msal iyi uygulamalarla uyumludur. Aksiyon da\u{g}{\i}l{\i}m{\i} grafi\u{g}i (\c{S}ekil~\ref{fig:full}e) bu bulguyu nicel olarak desteklemektedir.

\subsection{G\"orsel G\"osterim}
Pygame tabanl{\i} g\"orselle\c{s}tirme tarlay{\i}, bitki sa\u{g}l{\i}\u{g}{\i}n{\i} ve ajan{\i}n g\"unl\"uk kararlar{\i}n{\i} g\"ostermektedir. G\"ubre uygulamalar{\i} ve sulama olaylar{\i} hafif par\c{c}ac{\i}k efektleriyle belirtilmekte; \"ust bilgi ekran{\i} g\"un, tahmini verim ve se\c{c}ilen aksiyonu raporlamaktad{\i}r. Ortaya \c{c}{\i}kan animasyon s{\i}ral{\i} karar s\"urecini uzman olmayan izleyiciler i\c{c}in sezgisel hale getirmekte ve politikan{\i}n pratik sonu\c{c}lar{\i}n{\i} g\"orsel olarak do\u{g}rulamaktad{\i}r.

\section{Tart{\i}\c{s}ma}
Hibrit yap{\i}, taban verim ve g\"ubre istatistiklerinin ampirik kalibrasyonu ile \c{s}effaf stres tabanl{\i} b\"uy\"ume modelinin birle\c{s}imi olarak saf sim\"ulasyon ile veri k{\i}tl{\i}\u{g}{\i} aras{\i}nda pratik bir denge sunmaktad{\i}r. Stres katsay{\i}lar{\i} su ve besin stresi i\c{c}in literat\"ur aral{\i}klar{\i}ndan esinlendi\u{g}i i\c{c}in~\cite{doorenbos1979,raes2009aquacrop} dinamikler tar{\i}msal olarak yorumlanabilir kalmaktad{\i}r. Ayn{\i} zamanda model derin RL'nin gerektirdi\u{g}i y\"uksek b\"ol\"um say{\i}s{\i}n{\i} destekleyecek kadar hafiftir.

\"O\u{g}renilen politikan{\i}n sade g\"or\"unmesi (a\u{g}{\i}rl{\i}kl{\i} olarak hafif sulama ve 40\,kg g\"ubre) rastgele veya dengesiz bir sonu\c{c} de\u{g}ildir. Verilen ufuk ve stres dinamiklerinde ajan gereksiz y\"uksek dozlar{\i} elemi\c{s} ve sezonluk toplam{\i} ger\c{c}ek veri medyan{\i}na yakla\c{s}t{\i}rm{\i}\c{s}t{\i}r. Bu sadele\c{s}me, RL ajanlar{\i}n{\i}n s{\i}k\c{c}a sergiledi\u{g}i bir \"ozelliktir.

Bununla birlikte \c{c}al{\i}\c{s}man{\i}n s{\i}n{\i}rlar{\i} a\c{c}{\i}k\c{c}a belirtilmelidir. \"Oncelikle, istenen ideal veri seti mevcut de\u{g}ildir. {\.I}deal durumda her g\"un i\c{c}in nem, besin, uygulanan g\"ubre ve sulama miktar{\i} ile hasat sonucunu birlikte i\c{c}eren uzun bir zaman serisi kullan{\i}l{\i}rd{\i}. Eldeki veri seti ise b\"uy\"uk \"ol\c{c}\"ude sezonluk \"ozet d\"uzeyinde verim ve g\"ubre istatistikleri sa\u{g}lamaktad{\i}r; g\"unl\"uk karar--sonu\c{c} \c{c}iftleri do\u{g}rudan g\"ozlenememektedir. Bu nedenle ger\c{c}ek veri, ge\c{c}i\c{s} dinamiklerini sat{\i}r sat{\i}r \"o\u{g}retmek i\c{c}in de\u{g}il, taban verim ve g\"ubre da\u{g}{\i}l{\i}m{\i}n{\i} kalibre etmek i\c{c}in kullan{\i}lm{\i}\c{s}t{\i}r. G\"unl\"uk ge\c{c}i\c{s}ler stres tabanl{\i} model ile \"uretilmeye devam etmektedir. Bu, tamamen veri-g\"ud\"uml\"u bir \"o\u{g}renme iddias{\i}ndan farkl{\i}, hibrit bir yakla\c{s}{\i}md{\i}r.

{\.I}kinci olarak, problem modern derin peki\c{s}tirmeli \"o\u{g}renme \"ol\c{c}\"utlerine g\"ore g\"oreli olarak basittir. K{\i}sa ufuk, k\"u\c{c}\"uk ayr{\i}k aksiyon k\"umesi ve \c{s}ekillendirilmi\c{s} \"od\"ul, \"o\u{g}renmeyi kolayla\c{s}t{\i}rmaktad{\i}r. \c{C}al{\i}\c{s}man{\i}n katk{\i}s{\i} zor bir k{\i}yaslama g\"orevini k{\i}rmak de\u{g}il; tabular Q-\"O\u{g}renme'den DQN'e ge\c{c}i\c{s}, ger\c{c}ek istatistiklerle kalibrasyon ve rastgele politikaya kar\c{s}{\i} sistematik kar\c{s}{\i}la\c{s}t{\i}rmad{\i}r.

Di\u{g}er s{\i}n{\i}rl{\i}l{\i}klar aras{\i}nda sabit 30 g\"unl\"uk ufuk, tarla i\c{c}i mek\^ansal heterojenli\u{g}in olmamas{\i}, \c{c}evrimi\c{c}i hava tahminlerinin bulunmamas{\i} ve tam AquaCrop/DSSAT yerine sadele\c{s}tirilmi\c{s} b\"uy\"ume modeli say{\i}lmaktad{\i}r. Gelecek \c{c}al{\i}\c{s}malar daha uzun mevsimsel ufuklar, s\"urekli aksiyon uzaylar{\i} (PPO veya SAC), g\"unl\"uk \"ol\c{c}\"uml\"u tarla veri setleri ve s\"ure\c{c} tabanl{\i} sim\"ulat\"orlerle daha s{\i}k{\i} ba\u{g}lant{\i} i\c{c}erebilir~\cite{overweg2021cropgym}.

\section{Sonu\c{c}}
Bu \c{c}al{\i}\c{s}mada ak{\i}ll{\i} g\"unl\"uk g\"ubreleme ve sulama i\c{c}in DQN tabanl{\i} bir \c{c}\"oz\"um sunulmu\c{s}tur. Ger\c{c}ek \"ur\"un verim istatistikleri ile hafif fizik tabanl{\i} b\"uy\"ume modelinin birle\c{s}tirilmesi sayesinde ajan, toksisiteyi kontrol alt{\i}nda tutarken verimi rastgele politikaya g\"ore art{\i}ran politikalar \"o\u{g}renmi\c{s}tir. Sezonluk g\"ubre kullan{\i}m{\i} ger\c{c}ek veri medyan{\i}na yak{\i}n seyretmi\c{s}; e\u{g}itim tan{\i}lar{\i} ve g\"orsel kar\c{s}{\i}la\c{s}t{\i}rma politikan{\i}n katk{\i}s{\i}n{\i} desteklemi\c{s}tir.

Sonu\c{c}lar, ideal g\"unl\"uk tar{\i}msal zaman serisinin yoklu\u{g}u ve problemin DRL a\c{c}{\i}s{\i}ndan g\"oreli basitli\u{g}i nedeniyle s{\i}n{\i}rl{\i} bir kapsamda yorumlanmal{\i}d{\i}r. Katk{\i}, m\"ukemmel bir tarla otonomisi iddias{\i} de\u{g}il; tabular yakla\c{s}{\i}mdan DQN'e ge\c{c}i\c{s}in, hibrit kalibrasyonun ve tekrarlanabilir bir deney d\"uzene\u{g}inin ortaya konmas{\i}d{\i}r.

\begin{thebibliography}{00}
\bibitem{mnih2015human}
V.~Mnih vd., ``Human-level control through deep reinforcement learning,'' \textit{Nature}, c.~518, no.~7540, ss.~529--533, 2015.

\bibitem{raes2009aquacrop}
D.~Raes, P.~Steduto, T.~C.~Hsiao ve E.~Fereres, ``AquaCrop---The FAO crop model to simulate yield response to water: II. Main algorithms and software description,'' \textit{Agronomy Journal}, c.~101, no.~3, ss.~438--447, 2009.

\bibitem{hsiao2009aquacrop}
T.~C.~Hsiao vd., ``AquaCrop---The FAO crop model to simulate yield response to water: III. Parameterization and testing for maize,'' \textit{Agronomy Journal}, c.~101, no.~3, ss.~448--459, 2009.

\bibitem{steduto2009aquacrop}
P.~Steduto, T.~C.~Hsiao, D.~Raes ve E.~Fereres, ``AquaCrop---The FAO crop model to simulate yield response to water: I. Concepts and underlying principles,'' \textit{Agronomy Journal}, c.~101, no.~3, ss.~426--437, 2009.

\bibitem{doorenbos1979}
J.~Doorenbos ve A.~H.~Kassam, ``Yield response to water,'' FAO Irrigation and Drainage Paper No.~33, Food and Agriculture Organization of the United Nations, Roma, 1979.

\bibitem{overweg2021cropgym}
H.~Overweg, H.~N.~C.~Berghuijs ve I.~N.~Athanasiadis, ``CropGym: A reinforcement learning environment for crop management,'' arXiv:2104.04326, 2021.

\bibitem{wang2025fertilization}
Z.~Wang vd., ``Reinforcement learning-based agricultural fertilization and irrigation,'' \textit{Agriculture}, c.~15, no.~8, s.~252, 2025.

\bibitem{salvagiotti2009nitrogen}
F.~Salvagiotti vd., ``Nitrogen uptake, fixation and response to fertilizer N in soybeans: A review,'' \textit{Field Crops Research}, c.~108, no.~1, ss.~1--13, 2008.

\bibitem{cameron2013nitrogen}
K.~C.~Cameron, H.~J.~Di ve J.~L.~Moir, ``Nitrogen losses from the soil/plant system: A review,'' \textit{Annals of Applied Biology}, c.~162, no.~2, ss.~145--173, 2013.

\bibitem{tilman2002agricultural}
D.~Tilman vd., ``Agricultural sustainability and intensive production practices,'' \textit{Nature}, c.~418, ss.~671--677, 2002.

\bibitem{zhang2002precision}
N.~Zhang, M.~Wang ve N.~Wang, ``Precision agriculture---A worldwide overview,'' \textit{Computers and Electronics in Agriculture}, c.~36, no.~2--3, ss.~113--132, 2002.

\bibitem{madondo2023swat}
M.~Madondo vd., ``A SWAT-based reinforcement learning framework for crop management,'' AI for Social Good Workshop, 2023.

\bibitem{ba2016layer}
J.~L.~Ba, J.~R.~Kiros ve G.~E.~Hinton, ``Layer normalization,'' arXiv:1607.06450, 2016.

\bibitem{gautron2022gym}
R.~Gautron vd., ``gym-DSSAT: A crop model turned into a Reinforcement Learning environment,'' arXiv:2207.03270, 2022.

\bibitem{buet2022deep}
J.~Buet vd., ``Deep reinforcement learning for irrigation scheduling,'' \textit{Computers and Electronics in Agriculture}, 2022.

\bibitem{jones2003dssat}
J.~W.~Jones vd., ``The DSSAT cropping system model,'' \textit{European Journal of Agronomy}, c.~18, no.~3--4, ss.~235--265, 2003.

\bibitem{kamilaris2018deep}
A.~Kamilaris ve F.~X.~Prenafeta-Boldú, ``Deep learning in agriculture: A survey,'' \textit{Computers and Electronics in Agriculture}, c.~147, ss.~70--90, 2018.
\end{thebibliography}

\end{document}
